% !PRINTED 374
\weberpage{374}
6. \gesperrt{Die zum Führer $Q$ gehörige Ordnung $O$ heißt aus $O_1$ und $O_2$ zusammengesetzt oder komponiert.}

Wir setzen symbolisch
\weq{2}{O = O_1 O_2.}
Die Diskriminanten dieser drei Ordnungen sind
\weq{3}{D = Q^2 \Delta, \quad D_1 = Q_1^2 \Delta, \quad D_2 = Q_2^2 \Delta.}
Sind $Q_1$ und $Q_2$ relativ prim, so ist die aus beiden komponierte Ordnung die Hauptordnung $O_0$.

Um die Komposition der Ordnungen auszuführen, kann man ebenso verfahren, wie im vorigen Paragraphen. Wir setzen
\weq{4}{Q_1 = Qm_1, \quad Q_2 = Qm_2,}
und dann sind $m_1$, $m_2$ relativ prim zueinander.

Es mögen dann $\mathfrak{a}$ und $\mathfrak{b}$ zwei Ideale bedeuten, die zu $Q_1 Q_2$ relativ prim sind.

Ist dann
\[
\begin{aligned}
\mathfrak{a}&\ \text{äquivalent}\ \mathfrak{a}'&&\ \text{nach}\ O_1,\\
\mathfrak{b}&\ \text{\glqq}\ \mathfrak{b}'&&\ \text{\glqq}\ O_2,
\end{aligned}
\]
so ist
\[
\begin{aligned}
\frac{\mathfrak{a}}{\mathfrak{a}'}&=\eta_1\ \text{eine Zahl in}\ O_1,\\
\frac{\mathfrak{b}}{\mathfrak{b}'}&=\eta_2\ \text{\glqq}\ \text{\glqq}\ \text{\glqq}\ O_2,
\end{aligned}
\]
d.\,h. $\eta_1$ ist nach dem Modul $Q_1$, $\eta_2$ nach dem Modul $Q_2$ mit einer rationalen Zahl kongruent. Also sind beide, und folglich ihr Produkt, nach dem Modul $Q$ mit einer rationalen Zahl kongruent und gehören in die Ordnung $Q$. % [sic] Gedruckt ist Q statt O.
Demnach ist auch
\[
\mathfrak{c} = \mathfrak{a}\mathfrak{b} \ \text{äquivalent mit} \ \mathfrak{c}' = \mathfrak{a}'\mathfrak{b}' \ \text{nach} \ O.
\]
Es sei $a_1$ die kleinste durch $\mathfrak{a}$ teilbare natürliche Zahl und $\mathfrak{b}$ relativ prim zu $a_1$ (folglich auch zu $\mathfrak{a}$), und wenn $a_2$ die kleinste durch $b$ teilbare natürliche Zahl ist, so sind auch $a_1$ und $a_2$ relativ prim. % [sic] In der Quelle steht das zweite b hier in Antiqua; dem Zusammenhang nach ist das Ideal $\mathfrak{b}$ gemeint.
Man kann dann immer, was auch $b_1$, $b_2$ sein mögen, den Kongruenzen
\[
\begin{aligned}
b_1 &\equiv m_1 b \ (\mathrm{mod}\ 2a_1),\\
b_2 &\equiv m_2 b \ (\mathrm{mod}\ 2a_2)
\end{aligned}
\]
genügen und demnach die Basisformen von $\mathfrak{a}$ und $\mathfrak{b}$ in der Gestalt annehmen:
%CONT>>
